\documentclass[10pt,a4paper]{amsart}
\usepackage[margin=19mm]{geometry}
\usepackage{amsmath,amssymb,mathtools}
\usepackage[scr=rsfs,cal=euler]{mathalfa}
\usepackage{xfrac}
\usepackage[unicode,hidelinks,bookmarks=false]{hyperref}
\newcommand{\ie}{i.e.\ }
\newcommand{\cf}{cf.\ }
\newcommand{\resp}{respectively\ }
\newcommand{\Subsection}{\subsection}
\providecommand{\nobreakdash}{\nobreak\hbox{-}}
\setlength{\emergencystretch}{2em}
\multlinegap0pt
\usepackage{amssymb}
\newtheorem{proposition}{Proposition}
\newtheorem{definition}{Definition}
\newtheorem{lemma}{Lemma}
\newtheorem{theorem}{Theorem}
\newcommand{\fn}{\mathbb{E}\sp{1}}
\title{Best approximation results for fuzzy-number-valued continuous functions}
\author{Juan J. Font, Sergio Macario}
\date{Source: arXiv:2512.20667v1 (CC BY 4.0)}
\begin{document}
\maketitle

\noindent\emph{Source and licence.} Best approximation results for fuzzy-number-valued continuous functions. Version arXiv:2512.20667v1 (CC BY 4.0), distributed by arXiv under the Creative Commons Attribution 4.0 International licence, http://creativecommons.org/licenses/by/4.0/, as recorded in the arXiv licence metadata for this exact version and shown on the abstract page https://arxiv.org/abs/2512.20667v1. Full licence notice: \texttt{licenses/arxiv-2512.20667-CC-BY-4.0.txt}.\par\smallskip

\noindent\textit{Editorial scope.} Counted unit: the article's theorem main (source0.tex lines 415--419). The article's complete original proof of it is delivered with it (source0.tex lines 421--504, one \texttt{proof} environment of 84 lines). No mathematics is written, repaired or re-derived: the statement and the proof are the article's own lines, in the article's own notation, and no retained line is substituted or re-flowed.  Retained uncounted as the source-local context the proof needs: lines 272--279; lines 317--319; lines 321--328; lines 360--367.  Nothing was omitted from any retained span, so the package carries no omission record.  The retained lines cite 1 works, whose entries are transcribed verbatim from the article's own bibliography source in the e-print and delivered with short editorial labels recorded in the item's reference mapping.  No retained line contains a figure environment, an image-inclusion command or drawing code, so no image is delivered and the package declares no asset dependency.  Editorial formatting only, in addition: an AMS \texttt{amsart} preamble that restates the article's own declarations and macros used by the retained lines, namely the environments \texttt{proposition}, \texttt{definition}, \texttt{lemma}, \texttt{theorem} on separate counters, together with the standard packages those lines need.  The retained environments print in this document as Proposition 1, Definition 1, Lemma 1, Theorem 1 in order of appearance, whereas the article prints the same items with the numbering of its own full document.  The complete native source of the article as distributed in the arXiv e-print for this exact version is bundled unchanged as \texttt{sources/191587/source0.tex}.
\begin{proposition} \label{properties} (\cite[Proposition 2.3]{FSS17}) The metric space $(\fn, d_\infty)$ satisfies the following properties:
\begin{enumerate}
\item $d_\infty ( \sum_{i=1}^{m}u_{i}, \sum_{i=1}^{m}v_{i}) \le \sum_{i=1}^{m}d_{\infty} (u_i, v_i)$ where $u_i,v_i\in \mathbb{E}^1$ for $i=1,...,m$.
\item  $d_\infty ( ku, kv) = k d_\infty (u, v)$ where $u,v\in \mathbb{E}^1$ and $k>0$.
\item $d_\infty (ku,\mu u ) = \mid k - \mu  \mid d_\infty (u,0),$ where $u \in \mathbb{E}^1$, $ k \ge 0$ and $\mu \ge 0$.
\item $d_\infty (ku,\mu v ) \le \mid k - \mu  \mid d_\infty (u,0) + \mu d_\infty ( u,v),$ where $u,v \in \mathbb{E}^1$, $ k \ge 0$ and $\mu \ge 0$.
\end{enumerate}
\end{proposition}

\begin{definition} \label{multi} Let $W$ be a nonempty subset of $C(K, \mathbb{E}^1)$. We define
$$Conv(W):=\{ \varphi \in C(K, [0,1]): \varphi f + (1 - \varphi) g \in W \hspace{0,2cm} for \hspace{0,2cm} all \hspace{0,2cm} f,g \in W \}.$$
\end{definition}

\begin{proposition} \label{multiplier} (\cite[Proposition 3.2]{FSS17}) Let $W$ be a nonempty subset of $C(K, \mathbb{E}^1)$. Then we have:
\begin{enumerate}
\item $\phi \in Conv(W)$ implies that $1 - \phi \in Conv(W)$.
\item If $\phi, \varphi \in Conv(W)$, then $\phi \cdot \varphi \in Conv(W)$.
%\item If $\phi$ belongs to the uniform closure of $Conv(W)$, then so does $1 - \phi$.
%\item If $\phi, \varphi$ belong to the uniform closure of $Conv(W)$, then so does $\phi \cdot \varphi$.
\end{enumerate}
\end{proposition}

\begin{lemma}\label{pat3} (\cite[Lemma 3.6]{FSS17}) Let $ W \subseteq C(K, \mathbb{E}^1)$. If $Conv(W)$ separates the points of $K$, then, given $x_0 \in K$ and an open neighborhood $N$ of $x_0$, there exists a neighborhood $U$ of $x_0$ contained in $N$ such that for all $0 < \delta < \frac{1}{2}$, there is $\varphi \in Conv(W)$ such that

\begin{enumerate}
\item $\varphi (t) > 1 - \delta,$ for all $t \in U$;
\item $\varphi (t) < \delta,$ for all $t \notin N$.
\end{enumerate}

\end{lemma}

\begin{theorem} \label{main} Let $W$ be a subspace of $C(K,\mathbb{E}^1)$ and assume that $Conv(W)$ separates points.
For each $f \in C(K,\mathbb{E}^1)$, we have
\[d(f,W)=d_x(f,W)\]
for some $x\in K$.
\end{theorem}

\begin{proof}
We shall first show that $d(f,W)=\sup_{x\in K}d_x(f,W)$. It is apparent that $d(f,W)\ge \sup_{x\in K}d_x(f,W)$ since $d(f,W)\ge d_x(f,W)$ for each $x\in K$.
Let us prove that $d(f,W)\le \sup_{x\in K}d_x(f,W)$.

To this end, fix $\varepsilon > 0$. Given $x'\in K$, we can find $f_{x'}\in W$ such that $d_{\infty}(f(x'),f_{x'}(x'))<\sup_{x\in K}d_x(f,W)+\varepsilon$.
Let us define the following open neighborhood of $x'$:
$$N(x'):= \{ t \in K: d_\infty (f(t), f_{x'}(t)) < \sup_{x\in K}d_x(f,W)+\varepsilon \}.$$
Let  $U(x')$ be an open neighborhood of $x'$ satisfying the properties in Lemma \ref{pat3}.


Since $K$ is compact, there exist finitely many $\{ x_1,	\ldots , x_m \}$ in $K$ such that $K \subset U(x_1) \cup \ldots \cup U(x_m)$.
Choose $\delta>0$ such that $\delta<min(1, \frac{\varepsilon}{km})$, where
\[k:=\max \{D(f,0), D(f, f_{x_1}),...D(f, f_{x_m})\}.\]
By Lemma \ref{pat3}, we know that there exist $\phi_1, \cdots, \phi_m \in Conv(W)$ such that for all $i=1, \ldots, m$,
\begin{enumerate}
\item[(i)] $\phi_i (t) > 1 - \delta,$ for all $t \in U(x_i)$;
\item[(ii)] $0 \le \phi_i (t) < \delta,\quad if  \enspace  t \notin N(x_i)$.
\end{enumerate}

Let us define the functions
\newline
$\psi_1 := \phi_1$,
\newline
$\psi_2 := ( 1 - \phi_1) \phi_2$,
\newline
$\vdots$
\newline
$\psi_m := ( 1 - \phi_1)( 1 - \phi_2) \cdots ( 1 - \phi_{m-1}) \phi_{m}$.

By Proposition \ref{multiplier}, we know that $\psi_i \in Conv(W)$ for all $i=1,\ldots,m$.
Next we claim that
$$\psi_1+\ldots+\psi_j= 1 - ( 1 - \phi_1)( 1 - \phi_2) \cdots ( 1 - \phi_j),$$ $j=1,\ldots,m .$
Indeed, it is clear that
$$\psi_1 + \psi_2 = \phi_1  + (1 - \phi_1) \phi_2= 1-(1-\phi_1)\cdot(1-\phi_2).$$
We proceed by induction. Assume that the result is true for a certain $j\in\{4,...,m-1\}$ and let us check
$$\psi_1+\ldots+\psi_j+\psi_{j+1}= 1 - ( 1 - \phi_1)( 1 - \phi_2) \cdots ( 1 - \phi_j) ( 1 - \phi_{j+1}).$$
Namely,
$$\psi_1+\ldots+\psi_j+\psi_{j+1}=$$
$$=1 - ( 1 - \phi_1)( 1 - \phi_2) \cdots ( 1 - \phi_j)+( 1 - \phi_1)( 1 - \phi_2) \cdots ( 1 - \phi_{j}) \phi_{j+1}=$$
$$=1 - ( 1 - \phi_1)( 1 - \phi_2) \cdots ( 1 - \phi_j) ( 1 - \phi_{j+1}),$$
as was to be checked.
%Finally, we can define $\psi_1 := (1 - \phi_2) \cdots (1 - \phi_m)$, which belongs to $Conv(W) $ and satisfies $\psi_1 + \psi_2+ \ldots + \psi_m = 1$.

Fix $x_0\in K$. Then there is some $i_0\in \{1, \ldots, m\}$ such that $x_0\in U(x_{i_0})$. Hence $\phi_{i_0}(x_0)> 1 - \delta$ and consequently,
\[1\ge \sum_{i=1}^{m}\psi_i(x_0)=1-(1 - \phi_{i_0}(x_0)) \prod_{i \not= i_0} ( 1 - \phi_i(x_0))> 1- \delta.\]

Furthermore, we clearly infer
\begin{equation}\label{1}
\psi_i(t) < \delta \hspace{0.1in} {\rm for} \hspace{0.04in} {\rm all} \hspace{0.04in} t \notin N(x_{i_0}), \hspace{0.04in}i = 1, \ldots, m.
\end{equation}

Let
\begin{equation}\label{g}
h := \psi_1 f_{x_1} + \psi_2 f_{x_2} + \ldots + \psi_m f_{x_m}.
\end{equation}
It seems apparent that
$$h= \phi_1 f_{x_1} + ( 1 - \phi_1 ) [\phi_2 f_{x_2} + (1 - \phi_2) [\phi_3 f_{x_3} + \cdots + (1 - \phi_{m-1}) [\phi_m f_{x_m} \cdots]].$$
Therefore, $h \in W$ since $\phi_i\in Conv(W)$ for $i=1,...,m$ (see Definition \ref{multi}).

From Proposition \ref{properties}, we know that, given $x \in K$,
$$d_\infty (f(x), h(x)) \le d_\infty \left(f(x),\sum^m_{i=1} \psi_i (x)f(x)\right) + d_\infty \left(\sum^m_{i=1} \psi_i (x)f(x),h(x)\right)\le $$
$$\le \left|1-\sum^m_{i=1}\psi_i (x)\right|d_\infty \left(f(x), 0)\right) + \sum^m_{i=1}\psi_i (x) d_\infty \left( f(x),f_{x_i}(x)\right).$$

On the one hand, $\left|1-\sum^m_{i=1}\psi_i (x)\right|d_\infty \left(f(x), 0)\right)\le \delta D(f,0) \le \epsilon$.

On the other hand, let $$I_x=\{ 1 \le i \le m: x \in N(x_i) \}$$ and $$J_x=\{ 1 \le i \le m: x \notin N(x_i) \}.$$
Then, for all $i \in I_x$, we have
$$\psi_i (x) d_\infty (f(x), f_{x_i}(x)) \le \psi_i(x) (\sup_{x\in K}d_x(f,W)+\varepsilon)\le \sup_{x\in K}d_x(f,W)+\varepsilon$$
and, for all $i \in J$, the inequality (\ref{1}) yields
$$\psi_i(x) d_\infty (f(x),f_{x_i}(x)) \le \delta d_\infty (f(x),f_{x_i}(x))\le \delta D(f, f_{x_i})\le \delta k.$$

From the above two paragraphs, we can infer
$$d_\infty (f(x), h(x)) \le \epsilon + \sup_{x\in K}d_x(f,W)+\varepsilon + \delta k m \le \sup_{x\in K}d_x(f,W)+3\varepsilon$$
and, since $x\in K$ is arbitrary,
$$D(f,h) \le \sup_{x\in K}d_x(f,W)+3\varepsilon.$$
As a consequence, we deduce that
$$d(f,W)=\inf_{g\in W}\{D(f,g)\}\le \sup_{x\in K}d(f_x,W_x).$$

Finally, we can define a continuous function $\gamma: K\longrightarrow \mathbb{R}$ as
\[\gamma(x):=\inf_{g\in W}\{d_{\infty}(f(x),g(x))\}.\]
Since $K$ is compact, we know that $\gamma$ attains its supremum at some $x'\in K$. Hence, we can write
\[d(f,W)=d_{x'}(f,W).\]

\end{proof}

\begin{thebibliography}{99}
\bibitem[FSS17]{FSS17} J. J. Font, D. Sanchis, M. Sanchis, \textit{A version of the Stone-Weierstrass theorem in fuzzy analysis}. J. Nonlinear Sci. Appl. {\bf 10} (2017), 4275--4283. %
\end{thebibliography}
\end{document}
